\paragraph{Data.} Classification uses five public CT cohorts with \numLabelsN{} binary labels (Tabs.~\ref{tab:cohorts} and~\ref{tab:labelcounts} and Fig.~\ref{fig:examples} in Appendix~\ref{app:cohorts}). The CVPR26 cohort is the release of the CVPR 2026 3D CT foundation-model challenge, which contains AMOS \citep{ji2022amos} with 15 disease labels, COVID-CTset \citep{rahimzadeh2021covidctset} and LUNA25 nodule malignancy \citep{peeters2026luna25}. FDG-PET/CT \citep{gatidis2022fdgpetct} provides four diagnosis groups, which we predict from the CT alone. The remaining cohorts are RSNA 2022 cervical-spine fracture detection (RSNA22 C-spine; \citealp{lin2023rsnacsfdata}), RSNA 2023 abdominal trauma detection (RSNA23 trauma; \citealp{rudie2024ratic}) and RSNA 2019 intracranial hemorrhage (ICH) detection (RSNA19 ICH; \citealp{flanders2020rsnaich}). Segmentation uses AMOS22 (15 organs; \citealp{ji2022amos}), KiTS23 \citep{heller2023kits21}, LiTS \citep{bilic2023lits}, MSD Pancreas \citep{antonelli2022msd} and TotalSegmentator \citep{wasserthal2023totalsegmentator,wasserthal2026totalsegdata}, whose 117 structures we merge into 15 (Tab.~\ref{tab:totalseg}), for \numStructuresN{} structures in total. For the CVPR26 cohort and AMOS22, the official validation cases are divided into a selection and a held-out half, whereas the other sources use fixed seeded splits of the released labeled cases, and TotalSegmentator holds out its official test split.

\paragraph{Optimization.} All runs use AdamW with a 10\% linear warm-up, cosine decay to $10^{-6}$ and bf16 precision. A classification run trains for 8000 steps at batch size 1 on a masked, positive-weighted binary cross-entropy, and a segmentation run trains for 8000 steps of four crops on the sum of soft Dice and binary cross-entropy. For the two-stage methods, we split the budget in halves. Consequently, a joint model uses $1/K$ of the steps of the $K$ per-task models it replaces, except for the 15-structure joint segmentation models (AMOS22, TotalSegmentator), which train for 24\,000 steps. Each method has one default learning rate, set by what it trains: $10^{-3}$ for the classification linear probe and the segmentation decoder, $10^{-2}$ for the segmentation linear probe, $10^{-4}$ for the attentive probe and the adapters, and $2\times10^{-5}$ for encoder weights. None of these settings was tuned on our data (Tabs.~\ref{tab:optim} and~\ref{tab:optimseg}). Every 1000 steps, the model is evaluated on the selection split, and the checkpoint with the best macro AUROC or mean Dice is kept and its held-out metric reported.

\paragraph{Evaluation.} We report AUROC per label and Dice per structure and case, on each backbone's working grid \citep{maierhein2024metrics}. Methods are compared per label-mode pair with DeLong's test \citep{delong1988} for AUROC or a paired case bootstrap for Dice, both with Benjamini-Hochberg correction \citep{benjamini1995fdr}, and are ranked within each pair that has every method. TotalSegmentator per-structure models were trained for Lin, Dec, LoRA and Full FT only. The main comparisons are summarized at label level, where each label counts once with its two modes averaged, with 95\% intervals from a bootstrap over cohorts, labels and held-out cases (Appendix~\ref{app:stats}). As the bootstrap draws only four or five cohorts per backbone, which makes its intervals optimistic, and does not include the variation between training seeds, we also count the cohorts with a positive mean difference. We flag a method as not trained when the median standard deviation of its held-out probabilities is below \numSpreadMinValue{} (Appendix~\ref{app:stats}). Only SPECTRE Full FT was flagged (\numSpreadSpectreFullft{}, against at least \numSpreadOthermin{} for every other method). We report it in the tables but exclude it from ranks and from averages over methods. The joint models of Lin, LoRA, Attn-only FT and Full FT (classification) and of Dec, LoRA and Full FT (segmentation) were trained with two further seeds (Appendix~\ref{app:revision}).
